\section{Gap-\texorpdfstring{$0$}{0} inequalities}
\label{sec:gap-zero}

For a rational distance vector $d$ on $n$ vertices, let
$Z=\one\one\trans-2D$, where $D$ has off-diagonal entries $d_{ij}$ and
zero diagonal. Thus $Z_{ii}=1$. A gap inequality has the form
$b\trans Zb\ge\gamma(b)^2$, where
$\gamma(b)=\min_{s\in\{\pm1\}^n}|s\trans b|$.

\begin{lemma}\label{lem:gap-zero-hyperplanes}
The point $d$ violates a gap-$0$ inequality if and only if
$Z$ restricted to $s^\perp$ is not positive semidefinite for some
$s\in\{\pm1\}^n$. In particular, gap-$0$ inequalities never separate
a positive semidefinite $Z$.
\end{lemma}

\begin{proof}
An integer vector $b$ has $\gamma(b)=0$ exactly when $s\trans b=0$ for
some sign vector $s$. Hence a gap-$0$ violator gives a negative
direction on $s^\perp$. Conversely, fix such an $s$ and use the rational
basis with columns $e_i-s_is_ne_n$, $1\le i<n$, for $s^\perp$.
If the restricted form is not positive semidefinite,
\cref{lem:rational-negative-direction} gives a rational negative
direction, which can be scaled to an integer vector $b\in s^\perp$.
Then $b\trans Zb<0=\gamma(b)^2$.
\end{proof}

\begin{theorem}\label{thm:gap-zero-hardness}
Deciding whether a rational distance vector violates a gap-$0$ inequality is
$\NP$-complete. Hardness holds even for points in the interior of the
metric polytope whose matrix $Z$ has exactly one negative eigenvalue.
The reduction is from PARTITION and does not establish strong
$\NP$-hardness.
\end{theorem}

\begin{proof}
Membership in $\NP$ follows from \cref{lem:gap-zero-hyperplanes}: a sign
vector $s$ is a polynomial-size certificate, and exact elimination on
the displayed rational basis of $s^\perp$ verifies that the restricted
form is not positive semidefinite. It also supplies an integer violating
vector of polynomial encoding length.

We reduce PARTITION \citep{Karp1972}. Pad an instance by zeros so that
it contains $m\ge3$ numbers, and append another $m$ zeros. The resulting
vector $c\in\Z_{\ge0}^n$ has even length $n=2m\ge6$. A partition of the
original numbers can be balanced in cardinality by assigning signs to
the appended zeros; conversely, a balanced partition yields a partition
of the original numbers. We may restrict to
$C=\one\trans c>0$. The problem is therefore to find
$s\in\{\pm1\}^n$ with $s\trans c=s\trans\one=0$.

Set $K=4nC$, $a=c+K\one$, $A_2=\norm{a}^2$, and write $a_-$ and $a_+$
for the smallest and largest coordinates of $a$. Define
\begin{align*}
 \theta_*&=
 \frac{2(n-1)a_-^2-(n+1)a_+^2}{(n-1)(n-2)},
 &\nu&=\frac{\theta_*}{A_2a_+^2},\\
 \rho_i&=a_i^2\left(1+\frac{\nu a_i^2}{A_2}\right),
 &\Delta&=\sum_i\rho_i,\\
 \theta_{ij}&=\frac{\rho_i+\rho_j}{n-2}
              -\frac{\Delta}{(n-1)(n-2)}\quad(i\ne j).
\end{align*}
Let $L_\theta$ be the Laplacian with edge weights $\theta_{ij}$, let
$D_a=\diag(a)$, and set
\begin{equation}\label{eq:gap-zero-construction}
 Z=D_a^{-1}L_\theta D_a^{-1}-\nu\frac{aa\trans}{A_2}.
\end{equation}
All numbers in this construction are rational of polynomial encoding
length in $n+\operatorname{size}(c)$.

First we verify the weights. Since
$a_+/a_-\le1+1/(4n)$ and, for $n\ge4$,
$(n+1)(1+1/(4n))^2<2(n-1)$, we have $\theta_*>0$.
Moreover $\theta_*\le a_-^2/(n-1)$, so
\begin{equation}\label{eq:gap-zero-nu-bound}
 0<\nu\le\frac1{n(n-1)a_+^2}\le\frac1n.
\end{equation}
Writing $\rho_-$ and $\rho_+$ for the extreme $\rho_i$, we obtain
\[
 \theta_{ij}\ge
 \frac{2(n-1)\rho_- -n\rho_+}{(n-1)(n-2)}
 \ge\theta_*>0,
\]
because $\rho_-\ge a_-^2$ and
$\rho_+\le a_+^2(1+\nu)\le a_+^2(1+1/n)$.
A direct sum gives $\sum_{j\ne i}\theta_{ij}=\rho_i$.
Consequently $Z_{ii}=\rho_i/a_i^2-\nu a_i^2/A_2=1$, and
$Za=-\nu a$.

For $b\perp a$, put $h=D_a^{-1}b$. For every real $t$,
\[
 \sum_i a_i^2(h_i-t)^2=\norm{b}^2+t^2A_2\ge\norm{b}^2.
\]
Since every edge weight is at least $\theta_*$,
\[
 b\trans Zb=h\trans L_\theta h
 \ge n\theta_*\min_t\sum_i(h_i-t)^2
 \ge\frac{n\theta_*}{a_+^2}\norm{b}^2.
\]
Thus $Z$ is positive definite on $a^\perp$ and has exactly one negative
eigenvalue, namely $-\nu$.

If $s\trans a=0$, the integer vector $b=a$ has zero gap and
$a\trans Za=-\nu A_2<0$. Conversely, suppose
$s\trans b=0$ and $b\trans Zb<0$. Write $b=ta+b_\perp$, with
$b_\perp\perp a$. Since $Za=-\nu a$, the cross term vanishes, and
\[
 0> b\trans Zb
 \ge-\nu t^2A_2+
       \frac{n\theta_*}{a_+^2}\norm{b_\perp}^2.
\]
In particular $t\ne0$. Cauchy--Schwarz and $s\trans b=0$ give
\[
 |s\trans a|
 \le\frac{\sqrt n\norm{b_\perp}}{|t|}
 <\sqrt{\frac{\nu A_2a_+^2}{\theta_*}}=1.
\]
The quantity $s\trans a$ is integral, so it is zero. Finally,
$s\trans a=s\trans c+K s\trans\one$ vanishes exactly when both terms
vanish: $|s\trans c|\le C<K$ prevents cancellation of a nonzero
integer multiple of $K$. We have proved the required equivalence with
balanced PARTITION.

It remains to verify the metric-polytope promise. Put $r=a_+/a_-$.
For $n\ge4$, $r^2\le1+33/(64n)$ and $r^2<2$. As $a_+\ge K\ge4n$,
\eqref{eq:gap-zero-nu-bound} gives
\[
 \nu r^2\le\frac1{8n^3(n-1)}\le\frac1{384n},
 \qquad
 (1+\nu)r^2\le1+\frac{199}{384n}<1+\frac3{5n}.
\]
Using $\Delta\ge na_-^2$ and
$\rho_i\le(1+\nu)a_+^2$, we get, for $n\ge6$,
\[
 \frac{\theta_{ij}}{a_ia_j}
 \le\frac{2(n-1)(1+\nu)r^2-n}{(n-1)(n-2)}
 <\frac{n-4/5}{(n-1)(n-2)}\le\frac{13}{50}.
\]
The last inequality is equivalent to
$(n-6)(13n-11)\ge0$. Also
\[
 \frac{\nu a_ia_j}{A_2}
 \le\frac{\theta_*}{A_2^2}
 \le\frac1{16n^4(n-1)}<\frac1{100}.
\]
The cut coordinates from \eqref{eq:gap-zero-construction} are therefore
\[
 d_{ij}=\frac12+
         \frac{\theta_{ij}}{2a_ia_j}+
         \frac{\nu a_ia_j}{2A_2},
 \qquad \frac12<d_{ij}<\frac{127}{200}.
\]
They satisfy each ordinary triangle inequality strictly, since the
right-hand side exceeds $1$ and the left-hand side is below $1$.
Every triangle perimeter is less than $381/200<2$. Thus the point lies
in the interior of the metric polytope.
\end{proof}

The same promises hold arbitrarily close to the positive semidefinite
cone. In the construction one may replace $\nu$ by any smaller positive
rational value and recompute $\rho$ and $\theta$. The weight and metric
bounds remain valid, and the converse estimate becomes
$|s\trans a|<\sqrt{\nu A_2a_+^2/\theta_*}\le1$. Hence, for any prescribed
positive rational $\varepsilon$, the reduction can enforce
$0<\nu\le\varepsilon$ with polynomial encoding length in the source
instance and in $\varepsilon$. The distance to
the positive semidefinite cone in spectral norm is then at most
$\varepsilon$. Every gap-$0$ inequality is satisfied at positive
semidefinite points.

\begin{remark}\label{rem:general-gap-scope}
If $Z$ is not positive semidefinite, an integer negative direction gives
$b\trans Zb<0\le\gamma(b)^2$, so the decision version of unrestricted
gap separation is immediately positive. This also gives a polynomial-size
violating coefficient vector. Computing the exact gap of a prescribed
vector is a separate partition problem; the negative-direction argument
does not by itself compute the tight right-hand side of its gap
inequality. At positive semidefinite points, neither the gap-$0$ result
nor our reductions for hypermetric, gap-$1$ and rounded psd inequalities
settle separation of all gap inequalities. The reductions do not exclude
a gap violator with $\gamma(b)\ge3$ in a no-instance for rounded psd
separation.
\end{remark}
